\subsection{7. 相关工作}\label{related-work}

动态组合性与多个成熟的研究领域相交汇。我们综述其中最相关的若干条工作脉络，并区分我们的贡献与其中每一条的差异。

\subsection{7.1. 效应与余效应系统}\label{efect-and-coefect-systems}

第~2 节把效应与余效应作为支撑我们工作的理论基石作了回顾。我们首先定位如今在工业实践中常见的单子式效应系统，然后综述三条研究脉络——它们把效应与余效应朝与 Cordis 相关的方向作了扩展：把代数效应重新表述为能力、赋予效应可逆语义、以及在统一的分级纪律下把效应与余效应统一起来。

单子式效应系统。有一族库在既有通用语言的类型系统中对效应进行编码，把它们表示为由运行时执行的单子值。Scala 中的 ZIO {[}82{]} 把计算建模为 ZIO{[}R,E,A{]}，TypeScript 中的 Effect-TS {[}83{]} 则建模为 Effect\textless A,E,R\textgreater——一个泛型类型，其参数分别描述结果、带类型的错误以及上下文必须提供的服务；fp-ts 库 {[}84{]} 通过基于 Reader 的单子变换子来编码同样的错误通道与需求通道。有两点特征把这些系统与 Cordis 区分开来。其一，追踪是以单子式嵌入为代价换来的：程序只有写进效应类型内部才能获得这种追踪，而 Cordis 把效应追踪作为叠加在普通宿主代码之上的覆盖层。其二，需求通过解释来履行：由一个已安装的服务提供其操作，而该服务被撤走时，它的操作所执行的东西仍然原地留存；Cordis 则给每个效应配一个逆，并在提供者进出时重新解析需求（第~3.1 节、第~3.2 节）。

作为能力的代数效应。代数效应（第~2.1 节）使效应操作对类型系统可见。与我们的工作最接近的扩展是 Brachthäuser 等人的 Effekt 语言，它把效应类型重新解释为能力 {[}85, 86{]}：效应类型表达的是计算从其上下文要求什么，而不是它可能产生哪些副作用。这一视角与我们的相似，都把上下文视为能力的调解者。Cordis 与 Effekt 有两个方面的差异。(1) 在目的上，代数效应使效应可见是为了支持模块化解释，让一个操作拥有多种处理器语义；而 Cordis 使效应可见是为了支持追踪与逆转，给每一次上下文变换都配一个逆。(2) 在设定上，Effekt 在类型层面静态地约束效应，默认采用基于作用域的推理——能力是二等公民，被限制在其词法作用域内，并通过装箱（boxing）恢复头等用法，装箱通过在类型中追踪被捕获的能力来解除这一限制；Cordis 则在运行时约束效应，目标是组件移除时完整的资源恢复；第~6.7 节将讨论按此意义把上下文降为二等公民的语言会提供什么。

可逆效应语义。有一条平行脉络赋予效应可逆语义而非解释语义。Heunen 等人 {[}87{]} 通过把 Hughes 的箭头（arrow）改造为 dagger 箭头与逆箭头，在可逆设定中为副作用建模，从而刻画序列化、可变存储等操作具有逆的效应。这是与我们的可逆效应最接近的形式化论述：二者都主张给每个效应配上撤销它的手段，而不是通过处理器来履行它。二者的差异在于可逆性居于何处，以及它们对可逆性要求多少。Heunen 等人工作在指称性的范畴论框架中，其中可逆性是全局性质，由构造保证——因为每个计算都可逆，逆是双侧的，且从范畴结构中恢复出来。Cordis 在运行时追踪逆，而且对它们的要求更少：不是要求整个计算可逆，而是要求每个原子效应承认一个单侧逆，该逆由调用方在应用点提供而非推导出来，任何复合效应的逆由此按组合规则得出（第~3.1 节）。

作为统一效应与余效应的分级类型。Orchard 等人 {[}88{]} 提出把分级模态类型（graded modal types）作为总括性概念，同时涵盖效应推理（经由分级单子）与余效应推理（经由分级余单子），并在 Granule 语言中实现，表明单一类型系统既能追踪计算做了什么，也能追踪它需要什么；更近的工作把余效应扩展到类 Java 的命令式语言 {[}89, 90{]} 以及 call-by-push-value {[}91{]}。这些工作都在类型层面运作：效应与余效应是编译期在词法固定的作用域上检查的静态标注。我们的贡献与这一分析路线正交：我们把同样的两个概念提升为运行时机制，使 Cordis 能够处理动态组合。时间维撤回与空间维依赖会随已加载组件集合的演化而被重新解析，而不是在固定的程序文本上一次性地敲定。

\subsection{7.2. 编程范式}\label{programming-paradigms}

第~3.3.3 节确立了上下文范式，即一种通过显式上下文来介导效应与余效应的纪律。有两种既有范式值得作显式比较：一个与我们共享术语，另一个与我们对横切关注点的处理相通。

面向上下文编程。COP {[}92, 93{]} 为语言配备层（layer）——即部分的方法与类定义，它们按执行上下文在运行时被激活与停用，从而行为可以自适应，而基础代码无需指名其上下文依赖 {[}94{]}。COP 与 Cordis 都把上下文当作头等的、运行时可变的实体，并且都动态地激活与停用行为，但二者的相似只是名义上的。在 COP 中，``上下文'' 指环境性的执行情境（如位置、用户、模式），激活在某个动态作用域的范围内改变方法分派；层既不会追踪它引发的副作用，也不会逆转它们，激活亦不受依赖满足的约束。在 Cordis 中，上下文是介导效应与余效应的 \(\Gamma _ { \infty }\) 实体：激活运行组件的可逆效应，并由响应式余效应满足所驱动（第~3.2 节），停用则把它们完整逆转。COP 改变的是运行何种行为；Cordis 组合并逆转的是组件所安装的效应与依赖。二者的差异在于取舍不同。COP 把激活并入宿主语言的方法分派，以语言特异性为代价换来动态作用域的层扩展范围；而 Cordis 作为与语言无关的覆盖层，在共享上下文之上以响应式方式解析激活。因此 Cordis 只能把 COP 中全局的、由值驱动的那个片段表达为余效应：即实现之间的上下文相关选择，而非动态作用域的激活。

面向切面编程。AOP {[}95, 96{]} 把横切关注点模块化为切面：切点（pointcut）在基础程序选中的连接点（join point）上进行量化，通知（advice）则织入到每个连接点上。Cordis 处理的是同一个问题——否则会散布在各组件之间的上下文行为；但它的切面类似物是余效应：一个许多组件声明依赖的共同调解点，从而可以在不修改任何组件的情况下在那里重塑横切行为。两种范式随后在两个维度上分道扬镳。(1) 声明式对无感式（obliviousness）：AOP 的切点是无感的、量化的，匹配任意连接点，而连接点的代码并不知道自己被织入；Cordis 则把横切限定在各组件声明的余效应上，因此其触及范围恰好是所声明的表面。这带来确定性与可追溯性：应用编排器可以在配置层检查并管控组件被横切的内容，而无需阅读或分析其源码；AOP 关注点则只能通过对它进行量化的那些切面才可解读。(2) 生命周期集成：Cordis 中的横切变更由组件的效应承载，在组件卸载时被逆转，并响应式地传播给其依赖者，因此它是动态组合模型内部的一步操作；动态 AOP 系统 {[}97, 98{]} 也能在运行时织入与解织，但那是独立操作，既不受组件生命周期的约束，也不会在被织入的代码之间触发重新解析。

\subsection{7.3. 时间维组合性}\label{temporal-composability-1}

时间维组合性关注在运行中的程序里替换或移除组件，同时恢复它所安装的效应。既有方法按它们如何处理离场组件的状态与效应来划分：把状态前向携带给后继版本、通过开发者编写的清理逻辑恢复效应、在预先固定的作用域内自动逆转效应、或在运行时通过拦截某个接口所累积的记录来回收资源。

有状态前向迁移。有一大类系统通过把组件状态前向携带到新版本，从而在不停机的情况下替换运行程序中的组件。它们都遵循同一条时序纪律：组件只有在到达一个安全、无交互的点之后才能被换出。Kramer 与 Magee 把这一判据确立为静默（quiescence）{[}51{]}，Vandewoude 等人后来将其放宽为干扰更小的宁静（tranquility）{[}52{]}；我们的滚动更新模式（第~6.2 节）通过在卸载提供者之前排空在途请求来落实它。动态软件更新（dynamic software updating，DSU）随后通过手写的变换函数把状态前向迁移：Hicks 等人面向 C 语言的通用 DSU {[}99{]}、Stoyle 等人通过 con-freeness 分析得到的类型安全更新点 {[}100{]}、以及 Hayden 等人的 Kitsune {[}101{]}，都把旧版本数据映射到新版本表示，原地继承堆对象、打开的文件与连接，同时重新初始化那些未迁移的部分。同样的纪律也扩展到持久状态：Overeem 等人 {[}102{]} 通过手写的升级操作在架构版本之间转换运行中事件存储的数据，同时保持系统可用。Erlang/OTP {[}15{]} 在进程层面采取同样的立场，通过 code\_change/3 迁移状态，并通过重启被监督的进程来从故障中恢复，而不是逆转它们的效应；JavaScript 的热模块替换（hot module replacement，如 webpack {[}46{]}、Vite {[}47{]}）在模块层面如法炮制，在重载时通过 module.hot 或 import.meta.hot API 把状态前向交接。与 Cordis 的模块替换（第~5.2 节）相比，这些方法对内存中状态的迁移更优雅：Cordis 会逆转旧组件的已跟踪效应，并从干净的状态重新应用新组件的效应，因此组件自身的内存中状态在重载后不会留存，除非它被放在更长寿命的依赖中；在可逆效应之上叠加 DSU 式前向迁移属于未来工作。不过，Cordis 的方法在两方面更通用：它不需要 DSU 和 HMR 所需的那种手写迁移函数，并且它支持整体卸载组件并恢复其资源，而不仅仅是原地更新一个组件。

开发者编写的恢复。第二类方法通过开发者手工编写的清理或补偿逻辑来恢复组件的效应。插件生命周期约定（如 OSGi {[}50{]}、Eclipse 的扩展点、IntelliJ 与 VSCode）把清理委托给开发者编写的卸载回调；命令模式 {[}103{]} 把操作与撤销方法封装在一起，供撤销/重做栈使用；saga 模型 {[}49{]} 把长事务组织成若干步骤，每一步都配一个补偿动作；代数效应处理器可以在拆解（teardown）时附加并运行终结器 {[}104{]}；事件溯源 {[}105{]} 则通过追加补偿事件而非执行逆操作来撤回状态。在所有这些方法中，逆操作都是一种未被强制执行的义务，与操作本身相解耦，因此一旦忘记就会静默地泄漏资源（第~1.2.1 节对此有经验性的记录）。React 的 useEffect 钩子 {[}106{]} 在结构上最接近把效应与它的逆配对：它返回一个清理函数，运行时在每次重新执行前以及卸载（unmount）时调用该函数。它的不足在于组合性：钩子只能在组件或另一个钩子的顶层调用，绝不能出现在条件、循环或嵌套函数内，而且它的效应体既不能接受异步函数，也不能接受迭代器。因此效应无法由其他效应组装而成，也无法与控制流交错，也就无从派生出复合逆。Cordis 的效应没有这种限制：它们是可自由组合、可异步运行的普通操作，只需为每个原子效应手写一个逆，任何复合效应的逆都由组合规则派生出，因此组装已有效应完全不需要编写逆。这种把每个效应与其逆结构性地配对的做法，使完整恢复成为系统的不变量，而非开发者自律的问题。

静态作用域的逆转。第三类方法通过构造自动逆转效应，但把逆转限制在预先固定的作用域内。软件事务内存 {[}107, 108{]} 源自硬件事务内存 {[}109{]}，它记录读写日志，使一组内存操作要么提交要么中止，把内存回滚到事务前的状态。可逆计算从 Landauer 与 Bennett 的热力学分析 {[}110, 111{]} 到诸如 Janus {[}112{]} 之类的可逆语言，走得更远，使整个计算的每一步都全局可逆。可逆进程演算把回溯直接构建进语义本身：RCCS {[}113{]} 为每个进程携带一段记忆，当某一步引向的过去在因果上等价时，允许该步被撤回；Phillips 与 Ulidowski {[}114{]} 在保持 CCS、ACP、CSP 的前向操作语义的同时，统一地派生出它们的可逆算子。它们的因果一致性判据，正是 Cordis 恢复所遵循顺序的并发对应物：累加器以后进先出的顺序应用组件自身的逆，而第~4.3.1 节的守卫把提供者的撤出推迟到其消费者停用之后（定理~63）。然而，其触及范围由语义固定下来——所执行的每个动作都保持可撤销；而 Cordis 组件为每个原子效应提供一个逆，其累加器把上下文带回到组合开始之处。线性类型 {[}115{]}、RAII {[}4{]} 与 Rust 的所有权系统 {[}61{]} 把资源的释放绑定到词法区域。它们都在静态上固定逆转的作用域与触及范围；相比之下，Cordis 不预先固定这样的作用域：它在一个组件的生命周期内逆转任意的上下文操作，并把词法资源管理视为互补手段，适用于单个组件内的局部资源。

拦截式回收。第四类方法在组件自身不提供逆的情况下回收它所获取的资源，做法是在运行时控制的一个接口处记录它的获取行为。Nooks {[}77{]} 包装跨越 Linux 内核与其可加载扩展之间边界的每一次调用，使扩展所触及的内核对象都经过一个对象跟踪器，其记录告诉恢复管理器在扩展失败时该释放什么；影子驱动 {[}78{]} 从另一侧截取同样的调用，记录决定驱动状态的请求与配置，使重启后的实例得以恢复到该状态。Akeso {[}116{]} 改由编译器插桩来获得记录：它把内核执行划分为可嵌套的恢复域，记录域的状态变更与跨线程依赖，并把出错的请求连同依赖它的每个域一起回滚。因此回收源自运行时维护的记录，而非开发者记得去写的清理代码，这使这一类方法成为可逆效应在系统层面最接近的先例。它与 Cordis 的差别在于词汇与触及范围。平台决定什么可以被记录：或是按内核对象类型一份释放代码，或是每类驱动一个影子，或是每个被插桩的分配器一个逆——因此组件只能持有平台已经知道如何释放的资源；而 Cordis 组件会引入自己的效应，并为其中每个原子效应提供一个逆（第~3.1 节）。回收同样被一次提交的请求或同一扩展的重启所限制；而 Cordis 会在组件的整个生命周期内逆转，并把移除传播给其依赖者，由它们依次释放自己的效应（第~3.2 节）。

\subsection{7.4. 空间维组合性}\label{spatial-composability-1}

空间维组合性关注组件对其它组件的依赖如何声明与绑定。既有机制按绑定如何响应变化来划分：在初始化时一次性装配依赖、对整体组件的可用性作出反应、或在单个值的粒度上传播变化。

初始化时依赖装配。两种成熟的机制在初始化时把组件装配在一起。依赖注入框架 {[}38{]}（如 Spring {[}117{]}、Guice、Angular、Inversify）在初始化时把依赖注入组件，UI 框架的上下文（如 Vue.js 的 provide/inject 与 React 的 Context API）则沿组件树传递依赖。其中有些支持动态作用域（如 Spring 的 prototype/request 作用域、Angular 的分层注入器），但两者都不会响应式地重新解析：当提供者在运行时被替换或移除时，既有依赖者既不会被停用也不会被重新初始化，而且它们都不提供我们组件状态机所提供的那种生命周期管理。Cordis 的响应式余效应（第~3.2 节）正提供了这一点：每当满足谓词变化时，通知机制就触发生命周期转移。

可用性响应式组件模型。与我们响应式余效应最接近的先例对服务可用性作出反应。OSGi 的声明式服务（Declarative Services）与 iPOJO {[}118, 119{]} 允许组件声明所提供与所需的服务，运行时在服务出现与消失时自动激活与停用它们；iPOJO 的 Gravity 项目 {[}119{]} 明确以对不断变化的服务可用性进行自主运行时自适应为目标，其 provide/require 模型直接预示了 Cordis 的 ctx.provide/ctx.get 模式。R-OSGi {[}53{]} 通过 RPC 把同样的抽象透明地扩展到分布式环境，把网络故障映射为服务撤回事件——这是第~6.2 节作为 Cordis 模型的一种扩展来讨论的模式。所有这些系统都通过停用回调来恢复，而回调在两方面受限。其一，回调是手写的，因此资源安全依赖于开发者自律，一旦忘记就会静默泄漏。其二，回调是同步的：若拆解需要与即将离场的依赖进行异步交换，这些框架不提供任何等待它的协议，只能对可能已经过期的引用进行阻塞等待。Cordis 的响应式余效应填补了这两处空缺：停用会逆转依赖者累积的效应，而且其惯性 \(\mathsf{Unloading}\) 状态（第~4.3.3 节）会在对后续变化作出反应之前，把异步拆解完整地执行到底。

值级响应性。函数式响应编程（functional reactive programming，FRP）{[}120{]} 及其现代形态（如 SolidJS 中的信号 {[}121, 122{]}、Vue 的响应系统、Angular Signals）在值级粒度上传播变化：当信号变化时，派生计算被同步或在调度器下重新求值 {[}123{]}。Cordis 的响应式余效应在组件级粒度上运作，补充了值级传播所不建模的异步生命周期语义。同样的粒度差异在一致性上朝相反方向起作用：在依赖图固定顺序的一轮（turn）内传播，使 FRP 得以要求任何派生计算都不会读到新值与旧值混杂的输入，即无毛刺（glitch freedom）{[}124{]}；而 Cordis 没有对应的一轮概念，编排动作一个个到达，它只保证没有单个转移跨越两次余效应解析（定理~64）。二者互补而非竞争：Cordis 的余效应本身可以携带响应式值，组件只对其实际消费的部分进行更新，从而把组件级响应性细化为横跨两个层级的、粒度更细的响应式余效应。
